% !TEX TS-program = pdflatex
% MagnaDesign — Curso de uso interno (Beamer, sem deps externas)
%
% Build:
%   make           → produces main.pdf via pdflatex
%   make screenshots
%                  → regenerates figures/ from scripts/build_screenshots.py
%
% Why not metropolis? — keeps the deck portable across BasicTeX
% installs that don't ship the metropolis theme. The custom
% styling below replicates the metropolis aesthetic (sans-serif
% body, coloured frame titles, slim progress bar, dark accents)
% using only base beamer commands.

\documentclass[11pt, aspectratio=169]{beamer}

\usetheme{default}
\useinnertheme{default}
\useoutertheme{default}
\usefonttheme{professionalfonts}

\usepackage[T1]{fontenc}
\usepackage[utf8]{inputenc}
\usepackage[brazilian]{babel}
\usepackage{lmodern}      % Latin Modern — readable sans + clean
                          % at any zoom on screen
\usepackage{booktabs}
\usepackage{xcolor}
\usepackage{tikz}
\usetikzlibrary{positioning, arrows.meta, shapes.geometric, fit, backgrounds, calc}
\usepackage{listings}
\usepackage{graphicx}
\usepackage{xspace}

% MagnaDesign palette — synced with src/pfc_inductor/ui/theme.py.
% MagnaDesign palette — extracted from src/pfc_inductor/ui/theme.py.
% Keeping a single source of truth here so all custom blocks read
% from the same hex codes that the running app uses.
\definecolor{md-violet}{HTML}{8B5CF6}      % accent_violet
\definecolor{md-violet-dark}{HTML}{6D28D9} % accent_violet darker for emphasis
\definecolor{md-amber}{HTML}{F59E0B}       % warning
\definecolor{md-success}{HTML}{059669}     % success
\definecolor{md-danger}{HTML}{DC2626}      % danger
\definecolor{md-text}{HTML}{1F2937}        % text
\definecolor{md-muted}{HTML}{6B7280}       % text_muted
\definecolor{md-bg}{HTML}{F9FAFB}          % bg
\definecolor{md-surface}{HTML}{FFFFFF}     % surface
\definecolor{md-border}{HTML}{E5E7EB}      % border
\definecolor{md-code-bg}{HTML}{F3F4F6}     % code background

% Re-skin metropolis with our palette so the title bars / section
% accents pick up the app's identity. Metropolis exposes these
% style hooks for exactly this purpose.
\setbeamercolor{normal text}{fg=md-text, bg=md-bg}
\setbeamercolor{alerted text}{fg=md-danger}
\setbeamercolor{example text}{fg=md-success}
\setbeamercolor{frametitle}{fg=md-bg, bg=md-violet-dark}
\setbeamercolor{title separator}{fg=md-violet}
\setbeamercolor{progress bar}{fg=md-violet, bg=md-border}
\setbeamercolor{progress bar in head/foot}{fg=md-violet, bg=md-border}
\setbeamercolor{progress bar in section page}{fg=md-violet, bg=md-border}
\setbeamercolor{section in toc}{fg=md-violet-dark}
\setbeamercolor{block title}{bg=md-violet-dark, fg=md-bg}
\setbeamercolor{block body}{bg=md-code-bg, fg=md-text}
\setbeamercolor{block title alerted}{bg=md-danger, fg=md-bg}
\setbeamercolor{block title example}{bg=md-success, fg=md-bg}


% ─── Custom theme (metropolis-ish without the dependency) ───
%
% A minimalist look that survives any TeX install: sans-serif
% body, a thin coloured rule under the frame title, a discrete
% progress bar in the footer, and section pages with a coloured
% background. Each \setbeamer* below is independently optional;
% comment any out to dial back the styling.

% Fonts — Latin Modern Sans for everything (replaces Computer
% Modern Roman). Tighter line-height so each slide fits more
% concise content without wrapping.
\renewcommand{\familydefault}{\sfdefault}
\setbeamerfont{title}{size=\Large, series=\bfseries}
\setbeamerfont{subtitle}{size=\small, series=\mdseries, shape=\itshape}
\setbeamerfont{author}{size=\small}
\setbeamerfont{date}{size=\footnotesize}
\setbeamerfont{frametitle}{size=\large, series=\bfseries}
\setbeamerfont{framesubtitle}{size=\small, series=\mdseries, shape=\itshape}

% Frame title — coloured rule under it for the metropolis vibe.
\setbeamertemplate{frametitle}{%
  \nointerlineskip
  \begin{beamercolorbox}[wd=\paperwidth, leftskip=0.6cm,
                          rightskip=0.6cm, ht=2.4ex, dp=1ex]
    {frametitle}%
    \usebeamerfont{frametitle}\insertframetitle\strut\par
  \end{beamercolorbox}%
  {\color{md-violet}\rule{\paperwidth}{1.2pt}}%
}
\setbeamercolor{frametitle}{fg=md-text, bg=md-bg}

% Footer — small progress bar + section name on the left, slide
% number on the right.
\setbeamertemplate{footline}{%
  \leavevmode%
  \hbox{%
    \begin{beamercolorbox}[wd=.5\paperwidth, ht=2.5ex, dp=1ex,
                            leftskip=0.6cm]{author in head/foot}%
      \usebeamerfont{author in head/foot}%
      \color{md-muted}\insertshortauthor%
    \end{beamercolorbox}%
    \begin{beamercolorbox}[wd=.5\paperwidth, ht=2.5ex, dp=1ex,
                            rightskip=0.6cm]{title in head/foot}%
      \raggedleft\usebeamerfont{title in head/foot}%
      \color{md-muted}\insertsection\,\textbar\;%
      \insertframenumber\,/\,\inserttotalframenumber%
    \end{beamercolorbox}%
  }\vskip0pt%
}
\setbeamertemplate{navigation symbols}{}

% Section pages — coloured backdrop with the section name big.
\setbeamertemplate{section page}{%
  \begin{centering}
    \vskip 1cm
    {\color{md-violet}\rule{0.5\paperwidth}{1.4pt}}\par
    \vskip 0.7cm
    {\Large\bfseries\color{md-violet-dark}\insertsection}\par
    \vskip 0.5cm
    {\color{md-violet}\rule{0.5\paperwidth}{0.6pt}}\par
  \end{centering}
}
\AtBeginSection[]{%
  \begin{frame}[plain, noframenumbering]
    \sectionpage
  \end{frame}%
}

% Title page — clean, vertically centred.
\setbeamertemplate{title page}{%
  \vfill
  {\color{md-violet}\rule{0.6\paperwidth}{2pt}}\par
  \vskip 0.3cm
  {\usebeamerfont{title}\color{md-violet-dark}\inserttitle\par}
  \vskip 0.2cm
  {\usebeamerfont{subtitle}\color{md-text}\insertsubtitle\par}
  \vskip 1.0cm
  {\usebeamerfont{author}\color{md-muted}\insertauthor\par}
  {\usebeamerfont{date}\color{md-muted}\insertinstitute\,\textbar\;%
   \insertdate\par}
  \vfill
}

% Bullets — square accent dots in the brand colour.
\setbeamertemplate{itemize item}{\textcolor{md-violet}{\rule[0.4ex]{0.5em}{0.5em}}}
\setbeamertemplate{itemize subitem}{\textcolor{md-violet}{\rule[0.5ex]{0.4em}{0.4em}}}

% Block — coloured header strip + light fill body.
\setbeamertemplate{block begin}{%
  \begin{beamercolorbox}[colsep*=.75ex]{block title}%
    \usebeamerfont*{block title}\insertblocktitle\par
  \end{beamercolorbox}%
  \vskip-1pt%
  \begin{beamercolorbox}[colsep*=.75ex, vmode]{block body}%
    \usebeamerfont{block body}%
}
\setbeamertemplate{block end}{\end{beamercolorbox}}

% Standout-like frame — used for the closing "Q&A" slide.
\newenvironment{standoutpage}{%
  \begin{frame}[plain, noframenumbering]%
  \centering\vfill\Huge\bfseries\color{md-violet-dark}%
}{%
  \vfill\end{frame}%
}

% Code listings.
\lstdefinestyle{md}{
  basicstyle=\ttfamily\footnotesize\color{md-text},
  backgroundcolor=\color{md-code-bg},
  commentstyle=\color{md-muted}\itshape,
  keywordstyle=\color{md-violet-dark}\bfseries,
  stringstyle=\color{md-success},
  frame=none,
  breaklines=true,
  breakatwhitespace=true,
  showstringspaces=false,
  xleftmargin=1.2em,
  framexleftmargin=0.6em,
  numbers=none,
}
\lstset{style=md}

% Small helpers.
\newcommand{\app}{\textbf{MagnaDesign}\xspace}
\newcommand{\screenshot}[2][width=\linewidth]{%
  \begin{center}\includegraphics[#1]{#2}\end{center}}
\newcommand{\figcaption}[1]{{\footnotesize\color{md-muted}#1}}

\title{\app}
\subtitle{Engenharia de indutores PFC e reatores --- desktop tool}
\author{Curso interno}
\institute{Engenharia de Inversores de Frequência}
\date{\today}

\begin{document}

% ============================================================
%  ABERTURA
% ============================================================
\begin{frame}[plain, noframenumbering]
  \titlepage
\end{frame}

\begin{frame}{O que vamos cobrir hoje}
  \tableofcontents
\end{frame}

% ============================================================
\section{Visão geral do app}
% ============================================================

\begin{frame}{O problema que motivou a ferramenta}
  \begin{columns}[T,onlytextwidth]
    \column{0.55\textwidth}
    Engenharia de indutores PFC envolve hoje:
    \begin{itemize}
      \item \textbf{Várias planilhas} dispersas para cada topologia
      \item Validação \textbf{só por FEMM} (operador especializado)
      \item Difícil iterar sobre $L$, $N$, núcleo, fio sem refazer tudo
      \item Compliance EMC (IEC 61000-3-2) é cálculo separado
      \item Sem \textbf{rastro de design} entre projetos
    \end{itemize}
    \column{0.45\textwidth}
    \app integra:
    \begin{itemize}
      \item Engine analítico calibrado contra catálogos
      \item FEA automatizado (FEMMT/ONELAB + FEMM legacy)
      \item Otimizador multi-tier (cascade)
      \item Exportador estilo datasheet
      \item GUI única para 6 topologias
    \end{itemize}
  \end{columns}
\end{frame}

\begin{frame}{Topologias suportadas}
  \begin{table}
    \centering\small
    \begin{tabular}{@{}lll@{}}
      \toprule
      \textbf{Topologia} & \textbf{Faixa típica} & \textbf{Aplicações} \\
      \midrule
      Boost CCM & 200--2000 W & Server PSU, drives industriais \\
      Buck CCM & 100--1500 W & Conversores DC--DC industriais \\
      Interleaved Boost & 1--10 kW & PFC alta potência, ripple cancel.\\
      Flyback (DCM/CCM) & 30--300 W & SMPS, adaptadores, auxiliares \\
      Line reactor (3$\phi$) & 5--500 kW & VFD input, IEC 61000-3-2 \\
      Passive choke & 100 W--10 kW & PFC passivo, retificador 1$\phi$ \\
      \bottomrule
    \end{tabular}
  \end{table}
  \medskip

  Cada topologia tem o próprio cálculo analítico mas compartilha o
  mesmo motor de \textbf{perdas, térmica, FEA e otimização}.
\end{frame}

\begin{frame}{Arquitetura --- visão de blocos}
  \begin{center}
  \begin{tikzpicture}[
    every node/.style={font=\small},
    block/.style={draw=md-violet-dark, thick, rounded corners=2pt,
                  fill=white, minimum width=3cm, minimum height=0.8cm,
                  align=center},
    arrow/.style={-Stealth, md-violet-dark, thick},
  ]
    \node[block] (ui) at (0, 4)
      {GUI (PySide6)\\\scriptsize ProjetoPage / Análise / FEA dialog};
    \node[block] (eng) at (0, 2.5)
      {Engine analítico\\\scriptsize Numba JIT --- perdas + térmica};
    \node[block] (top) at (-3.2, 2.5)
      {Topologias\\\scriptsize Boost / Flyback / Reactor / ...};
    \node[block] (cat) at (3.2, 2.5)
      {Catálogo\\\scriptsize Cores + materiais + fios};
    \node[block] (fea) at (-3.2, 1.0)
      {FEA\\\scriptsize FEMMT / FEMM legacy};
    \node[block] (opt) at (0, 1.0)
      {Otimizador\\\scriptsize Pareto + cascade};
    \node[block] (exp) at (3.2, 1.0)
      {Export\\\scriptsize HTML / PDF datasheet};

    \draw[arrow] (ui) -- (eng);
    \draw[arrow] (eng) -- (top);
    \draw[arrow] (eng) -- (cat);
    \draw[arrow] (eng) -- (fea);
    \draw[arrow] (eng) -- (opt);
    \draw[arrow] (eng) -- (exp);
  \end{tikzpicture}
  \end{center}
  \medskip

  \figcaption{Engine no centro: cada superfície da UI fala com ele
  via Pydantic models (\texttt{Spec}, \texttt{Core}, \texttt{Wire},
  \texttt{Material}, \texttt{DesignResult}).}
\end{frame}

\begin{frame}{Tela principal --- visão da Análise}
  \screenshot{figures/example1_analise_overview.png}
  \figcaption{Núcleo + 3D + formas de onda + perdas + bobinamento +
  entreferro --- todos populando do mesmo \texttt{DesignResult}.}
\end{frame}

% ============================================================
\section{Exemplo 1 -- Boost PFC 1.5 kW}
% ============================================================

\begin{frame}{Exemplo 1 --- Boost PFC universal-input 1.5 kW}
  \begin{columns}[T,onlytextwidth]
    \column{0.5\textwidth}
    \textbf{Especificação}
    \begin{itemize}
      \item $V_\text{in} = 85$--$265$ V$_\text{rms}$
      \item $V_\text{out} = 400$ V (PFC bus)
      \item $P_\text{out} = 1500$ W
      \item $f_\text{sw} = 100$ kHz
      \item Ripple alvo: 30 \% de $I_\text{pk}$
    \end{itemize}

    \textbf{Referência de mercado}
    \begin{itemize}
      \item TI UCC28019 application note
      \item Magnetics Kool-Mu 0077439A7 toroide
      \item AWG 16 solid, 55--70 turns
    \end{itemize}

    \column{0.5\textwidth}
    \textbf{Por que esse exemplo?}
    \begin{itemize}
      \item Caso \textbf{mais comum} no mercado --- server PSU,
            telecom, drives industriais
      \item Bate com a janela de potência onde litz vs solid
            ainda é decisão de projeto
      \item Demonstra \textbf{rolloff de $\mu$} sob bias DC
            (Kool-Mu a 50--70\% no peak)
      \item Ativa todas as 5 abas de FEA
    \end{itemize}
  \end{columns}
\end{frame}

\begin{frame}{1.1 --- Spec drawer (entrada do usuário)}
  \screenshot{figures/example1_spec.png}
  \figcaption{O \texttt{SpecDrawer} valida cada campo via Pydantic e
  já dispara o engine ao apertar \texttt{Enter}. Topologia +
  faixa $V_\text{in}$ + $P_\text{out}$ + $f_\text{sw}$ ---
  3 cliques.}
\end{frame}

\begin{frame}{1.2 --- Análise: formas de onda + perdas}
  \begin{columns}[T,onlytextwidth]
    \column{0.6\textwidth}
    \screenshot{figures/example1_formas_onda.png}
    \column{0.4\textwidth}
    \textbf{Lê em segundos:}
    \begin{itemize}
      \item $i_L(t)$ no meio cycle de linha
      \item $B(t)$ em mT, comparado com $B_\text{sat}$
      \item Ripple peak-to-peak destacado
    \end{itemize}
    \figcaption{$P_\text{Cu,DC} \approx$ 1.1 W, $P_\text{Cu,AC}
    \approx$ 0.55 W, $P_\text{core} \approx$ 1.5 W. Total
    $\sim$ 3.15 W $\Rightarrow$ $\eta = 99.79\%$.}
  \end{columns}
\end{frame}

\begin{frame}{1.3 --- Loss stacked bar (proporção)}
  \screenshot[width=0.85\linewidth]{figures/example1_loss_stacked.png}
  \figcaption{Card de barra empilhada --- mostra qual família de
  perda domina. Cu DC concentra 35\%; revisar wire gauge vs Cu AC
  (skin/proximity) é o próximo passo.}
\end{frame}

\begin{frame}{1.4 --- Curva B-H operacional}
  \screenshot[width=0.78\linewidth]{figures/example1_bh_curve.png}
  \figcaption{Operating point dot (verde) na curva estática $B(H)$;
  linhas pontilhadas em $B_\text{sat}$ e a faixa vermelha de
  $-20\%$ headroom. Margin reportada: 66\% --- design conservador.}
\end{frame}

\begin{frame}{1.5 --- Geometria (vista datasheet)}
  \screenshot[width=0.65\linewidth]{figures/example1_geometry.png}
  \figcaption{Cross-section top-down do toroide --- 55 turns AWG16
  na bobbin radius, OD = 39.9 mm, ID = 24.0 mm. Datasheet stamp
  embaixo: \texttt{0077439A7 \textbullet{} 60 KoolMu \textbullet{} N = 55}.}
\end{frame}

\begin{frame}{1.6 --- Validação FEA: aba Summary}
  \screenshot{figures/example1_fea_summary.png}
  \figcaption{Bar chart pareado \emph{analítico vs FEA} para $L$ e
  $B_\text{pk}$. Confidence gauge embaixo (verde / amber / red)
  com pointer no pior dos dois erros. Modelo analítico
  \textbf{calibrado}.}
\end{frame}

\begin{frame}{1.7 --- FEA: heatmap do $|B|$ + 1D centerline}
  \begin{columns}[T,onlytextwidth]
    \column{0.5\textwidth}
    \screenshot{figures/example1_fea_field_heatmap.png}
    \begin{center}\scriptsize Cross-section r-z, viridis\end{center}
    \column{0.5\textwidth}
    \screenshot{figures/example1_fea_field_centerline.png}
    \begin{center}\scriptsize $|B|$ ao longo do gap (z = 0)\end{center}
  \end{columns}
  \figcaption{Cards do gallery --- clique abre lightbox com
  explicação por categoria.}
\end{frame}

\begin{frame}{1.8 --- FEA: $L$ vs corrente (rolloff)}
  \screenshot[width=0.85\linewidth]{figures/example1_fea_swept.png}
  \figcaption{Tier-4 swept-FEA roda solver em N pontos do bias.
  Linha violeta: $L(I)$. Tracejada: $B(I)$. Banda vermelha:
  $B_\text{sat}$. Triângulo: \emph{knee at $-30\%$} de $L$.}
\end{frame}

% ============================================================
\section{Exemplo 2 -- Reator de linha 3$\phi$ 22 kW}
% ============================================================

\begin{frame}{Exemplo 2 --- Reator de linha 3$\phi$ para VFD}
  \begin{columns}[T,onlytextwidth]
    \column{0.5\textwidth}
    \textbf{Especificação}
    \begin{itemize}
      \item Rede: 400 V$_\text{LL}$, 60 Hz, 3 fases
      \item Carga: 22 kW (drive de motor)
      \item $I_\text{rms} \approx 32$ A por fase
      \item Impedância alvo: \textbf{3 \%}
      \item $L \approx 2.5$ mH por fase
    \end{itemize}

    \textbf{Referência}
    \begin{itemize}
      \item ABB MCB-32 / Schaffner FN3220
      \item Núcleo ferro-silício M19
      \item $\sim$ 250 turns AWG 12 solid
      \item IEC 61000-3-2 Class A
    \end{itemize}

    \column{0.5\textwidth}
    \textbf{Por que esse exemplo?}
    \begin{itemize}
      \item Topologia \textbf{line-frequency} (60 Hz, sem
            comutação alta)
      \item Ativa o card de \textbf{harmônicas IEC 61000-3-2}
      \item N = 250 turns $\Rightarrow$ valida o auto-fallback
            FEMMT $\to$ FEMM legacy
      \item Sai na conformidade EMC ou não?
            Decisão visual em segundos
    \end{itemize}
  \end{columns}
\end{frame}

\begin{frame}{2.1 --- Spec do reator + análise inicial}
  \screenshot{figures/example2_spec.png}
  \figcaption{Topologia \texttt{line\_reactor} + \texttt{n\_phases =
  3} altera o cálculo de $L$ requerido (Erickson Cap. 18) e
  abre os campos específicos do reator (impedância, queda de
  tensão).}
\end{frame}

\begin{frame}{2.2 --- Card de harmônicas + IEC 61000-3-2}
  \screenshot[width=0.85\linewidth]{figures/example2_harmonics.png}
  \figcaption{Spectrum por ordem (h = 5, 7, 11, 13, 15) com limite
  IEC sobreposto como segmento vermelho. Bars verdes = abaixo
  do limite, amber = at margin, vermelho = falha. Verdict no
  título: \emph{PASSED} ou \emph{FAILED on h = X}. Triplens
  cancelam por simetria 3$\phi$.}
\end{frame}

\begin{frame}{2.3 --- Comparação antes/depois}
  \begin{columns}[T,onlytextwidth]
    \column{0.5\textwidth}
    \screenshot{figures/example2_harmonics_no_choke.png}
    \begin{center}\scriptsize Sem reator\end{center}
    \column{0.5\textwidth}
    \screenshot{figures/example2_harmonics_with_choke.png}
    \begin{center}\scriptsize Com 3\% reator\end{center}
  \end{columns}
  \figcaption{O $i(t)$ vira quase senoidal; a potência aparente
  (kVA) cai $\sim$ 22 \% para a mesma kW útil.}
\end{frame}

\begin{frame}{2.4 --- FEA do reator (N = 250 turns)}
  \begin{block}{Auto-fallback FEMMT $\to$ FEMM legacy}
    $N = 250$ excede o teto FEMMT (150 turns). O orchestrator
    \textbf{detecta automaticamente} e re-dispatcha para FEMM
    legacy (toroide axisymmetric, N sem custo geométrico).
    Usuário não vê erro --- só o resultado.
  \end{block}
  \medskip
  \figcaption{O ceiling existe porque FEMMT modela cada turn como
  curve loop separado em gmsh, que crashava em N alto. FEMM
  legacy modela o winding como bulk-current region: zero custo
  geométrico de $N$.}
\end{frame}

% ============================================================
\section{Exemplo 3 -- Flyback DCM 65 W}
% ============================================================

\begin{frame}{Exemplo 3 --- Flyback DCM 65 W}
  \begin{columns}[T,onlytextwidth]
    \column{0.5\textwidth}
    \textbf{Especificação}
    \begin{itemize}
      \item $V_\text{in} = 85$--$265$ V$_\text{rms}$
      \item $V_\text{out} = 19$ V (laptop)
      \item $I_\text{out} = 3.4$ A $\Rightarrow$ 65 W
      \item $f_\text{sw} = 65$ kHz, modo DCM
      \item $L_\text{p} \approx 350~\mu$H
    \end{itemize}

    \textbf{Referência}
    \begin{itemize}
      \item TI UCC28911 EVM (laptop adapter)
      \item Núcleo PQ20/16 ferrite N97
      \item $N_\text{p} = 42$, $N_\text{s} = 12$ ($n = 3.5$)
      \item Coupled inductor, 2-winding
    \end{itemize}

    \column{0.5\textwidth}
    \textbf{Por que esse exemplo?}
    \begin{itemize}
      \item Coupled-pair --- totalmente diferente
            do single-winding choke
      \item Ativa o render de \textbf{2 windings}
            (primary + secondary com cores distintas)
      \item Demonstra cálculo de \textbf{vazamento $L_\text{lk}$}
            (Mohan/Erickson formula)
      \item DCM --- current waveform triangular descontínuo
    \end{itemize}
  \end{columns}
\end{frame}

\begin{frame}{3.1 --- Spec do flyback}
  \screenshot{figures/example3_spec.png}
  \figcaption{Campos extras vs. boost: \texttt{Vout},
  \texttt{n\_turns ratio}, modo CCM/DCM inferido pela duty
  $D < 0.5$.}
\end{frame}

\begin{frame}{3.2 --- Geometry: cross-section coupled-pair}
  \screenshot[width=0.65\linewidth]{figures/example3_geometry.png}
  \figcaption{PQ20/16 com \textbf{primary em violeta} (42 turns,
  metade inferior da janela) e \textbf{secondary em âmbar}
  (12 turns, metade superior). Linha pontilhada divide as
  metades.}
\end{frame}

\begin{frame}{3.3 --- Análise + perdas do flyback}
  \begin{columns}[T,onlytextwidth]
    \column{0.55\textwidth}
    \screenshot{figures/example3_formas_onda.png}
    \column{0.45\textwidth}
    \begin{itemize}
      \item $i_\text{p}(t)$: triangular DCM, peak $\approx$ 4.2 A
      \item $i_\text{s}(t)$: refletida, peak $\approx$ 14.7 A
      \item Tempo morto evidente (DCM)
      \item $L_\text{lk} \approx 1.2~\mu$H (Mohan formula)
    \end{itemize}
    \figcaption{$V_\text{DS,pk} \approx 470$ V para FET 600 V padrão.}
  \end{columns}
\end{frame}

\begin{frame}{3.4 --- FEA validation flyback}
  \begin{columns}[T,onlytextwidth]
    \column{0.5\textwidth}
    \screenshot{figures/example3_fea_summary.png}
    \begin{center}\scriptsize Summary tab\end{center}
    \column{0.5\textwidth}
    \screenshot{figures/example3_fea_bh.png}
    \begin{center}\scriptsize B-H operacional\end{center}
  \end{columns}
  \figcaption{Para coupled inductors o FEA valida o $L_\text{p}$
  primário; o coupling $k$ é então derivado analiticamente do
  $L_\text{lk}/L_\text{p}$.}
\end{frame}

% ============================================================
\section{Recursos avançados}
% ============================================================

\begin{frame}{Otimizador Pareto (single-tier)}
  \screenshot[width=0.85\linewidth]{figures/feature_otimizador_pareto.png}
  \figcaption{Sweep cartesiano sobre (\texttt{core}, \texttt{material},
  \texttt{wire}, \texttt{N\_turns}) --- gera $\sim$ 1000--10000
  candidatos. Elimina dominados em (perdas, $\Delta T$, custo).
  Usuário inspeciona os 5--10 não-dominados.}
\end{frame}

\begin{frame}{Otimizador Cascade (multi-tier)}
  \begin{columns}[T,onlytextwidth]
    \column{0.55\textwidth}
    \begin{itemize}
      \item \textbf{Tier 1} --- analítico puro (segundos)
      \item \textbf{Tier 2} --- engine completo com skin/
            proximity ajustados (minutos)
      \item \textbf{Tier 3} --- FEA single-point no peak
            ($\sim$ 30 s/candidato)
      \item \textbf{Tier 4} --- swept-FEA $L(I)$ com rolloff
            (vários min/candidato)
    \end{itemize}
    \column{0.45\textwidth}
    \screenshot{figures/feature_cascade.png}
    \figcaption{Top-N table mostra valores refinados de cada tier.
    Coluna virtual $\text{COALESCE}(T4, T3, T2, T1)$ --- usuário
    sempre vê o número mais recente.}
  \end{columns}
\end{frame}

\begin{frame}{FEA Validation dialog --- 5 abas}
  \screenshot{figures/feature_fea_dialog_full.png}
  \figcaption{Summary (analítico vs FEA) | Geometry | B-H curve |
  Field plots (heatmap + 1D + histograma) | L vs current swept.}
\end{frame}

\begin{frame}{Cards específicos por topologia}
  \begin{columns}[T,onlytextwidth]
    \column{0.5\textwidth}
    \screenshot{figures/feature_phase_overlay_card.png}
    \begin{center}\scriptsize Phase overlay (interleaved)\end{center}
    \column{0.5\textwidth}
    \screenshot{figures/feature_harmonic_card.png}
    \begin{center}\scriptsize Harmonics (line reactor / passive)\end{center}
  \end{columns}
  \figcaption{Cards self-hide via \texttt{setVisible(False)} quando
  a topologia não casa --- a row colapsa e o usuário só vê o
  que importa pro design dele.}
\end{frame}

\begin{frame}{Compare designs (lado a lado)}
  \screenshot[width=0.85\linewidth]{figures/feature_compare.png}
  \figcaption{Tabela diff entre 2--3 \texttt{DesignResult}s. Highlight
  por tile de mudança. Útil pra justificar troca de núcleo ou
  material num design review.}
\end{frame}

\begin{frame}{Export --- datasheet HTML / PDF}
  \begin{columns}[T,onlytextwidth]
    \column{0.5\textwidth}
    \screenshot{figures/feature_export.png}
    \column{0.5\textwidth}
    \begin{itemize}
      \item Layout TDK / Würth / Vishay style
      \item Vista frontal/topo/lateral/iso (3D ortográfico)
      \item Tabela completa de specs
      \item Sessão de performance específica por topologia
      \item HTML self-contained (offline) ou PDF
    \end{itemize}
  \end{columns}
\end{frame}

\begin{frame}{Visualização 3D}
  \screenshot[width=0.6\linewidth]{figures/feature_3d.png}
  \figcaption{Render PyQt + Qt3D do toroide / E-core com winding
  visualizado. Útil pra confirmar montagem antes de mandar pro
  fornecedor de bobinamento.}
\end{frame}

% ============================================================
\section{Estado atual + roadmap}
% ============================================================

\begin{frame}{Estado atual}
  \begin{columns}[T,onlytextwidth]
    \column{0.5\textwidth}
    \textbf{Pronto e validado}
    \begin{itemize}
      \item 6 topologias com cálculo analítico
      \item FEMMT + FEMM legacy backends
      \item Otimizador single + cascade (4 tiers)
      \item Cards específicos por topologia
      \item Datasheet exporter (HTML)
      \item 5-aba FEA validation
    \end{itemize}
    \column{0.5\textwidth}
    \textbf{Performance}
    \begin{itemize}
      \item Engine: \textbf{3.6$\times$ speedup} via Numba
            kernel fundido
      \item Tier-1 analítico: $\sim$ 1 ms
      \item Tier-3 FEA point: 30 s
      \item Tier-4 swept (5 pts): 2--3 min
      \item Auto-update Sparkle-style + Ed25519 sig
    \end{itemize}
  \end{columns}
\end{frame}

\begin{frame}{Roadmap próximas releases}
  \begin{itemize}
    \item \textbf{Eficiência vs $P_\text{out}$} sweep ---
          validação cross-load (10--100\% de carga)
    \item \textbf{Voltage stress} per-component (FET, diodo)
    \item \textbf{Thermal coupled FEA} (FEMMT thermal solve)
    \item \textbf{Histórico de projetos} --- git-like timeline
          de iterações com diff visual
    \item \textbf{Magnetics catalog auto-update} via API do
          fornecedor (Magnetics, Micrometals, TDK)
  \end{itemize}
\end{frame}

% ============================================================
\section*{Encerramento}
% ============================================================

\begin{standoutpage}
  Perguntas?
  \\\medskip
  {\large\mdseries\color{md-muted}\app --- \today}
\end{standoutpage}

\begin{frame}{Referências bibliográficas}
  \footnotesize
  \begin{itemize}
    \item Erickson \& Maksimovi\'{c}, \emph{Fundamentals of
          Power Electronics}, 3rd ed., Cap. 18 (PFC).
    \item Mohan, \emph{Power Electronics}, 3rd ed., Cap. 4
          (flyback leakage).
    \item Hwu \& Yau, \emph{Resonance reduction in
          interleaved boost PFC}, IEEE TPE 2009.
    \item IEC 61000-3-2 Ed. 5 --- Harmonic emission limits.
    \item Magnetics Inc., \emph{Powder cores design manual}
          (Kool-Mu, High-Flux).
    \item Würth Elektronik, \emph{Trilogy of Magnetics},
          5th ed.
  \end{itemize}
\end{frame}

\end{document}
